\section{The Fundamental Trade-Offs: A Pareto Analysis}
\label{sec:tradeoffs}

\subsection{Four Evaluation Dimensions}

We evaluate each system on five dimensions that together capture the principal
objectives in the redistricting literature:

\begin{enumerate}
  \item \textbf{Compactness} --- measured by mean Polsby--Popper ratio
        $\mathrm{PP} = 4\pi A / P^2$ across all districts.  Higher is better;
        the baseline DIA achieves mean PP of 0.367 \citep{dellaLibera2026mec}.
        Compactness correlates with geographic coherence and resistance to
        partisan manipulation \citep{polsby1991third}.

  \item \textbf{Proportionality} --- measured as the absolute deviation of
        seat share from vote share across the 50 states, averaged by
        congressional delegation size.  Lower deviation is better.  The
        efficiency gap \citep{stephanopoulos2015partisan} provides a related
        but distinct measure; we report both where available.

  \item \textbf{Minority representation} --- measured as the number of
        majority-minority districts (MMDs) relative to the VRA Section~2
        benchmark derived from \emph{Thornburg v.\ Gingles}
        \citep{gingles1986} and interpreted through \emph{Allen v.\ Milligan}
        \citep{allen2023}.  We report the surplus (actual minus benchmark)
        to indicate whether a system over- or under-performs the legal
        requirement.\footnote{VRA score~3 = achieves majority-minority district
        targets in all covered states where minority population exceeds the 42\%
        threshold operationalised in D.1 \citep{dellaLibera2026d1}; score~2 =
        achieves targets in $\geq$3/5 covered states; score~1 = achieves
        targets in fewer than 3 covered states.  A majority-minority district
        is defined as a district in which a single minority group constitutes
        more than 50\% of the citizen voting-age population (CVAP), consistent
        with the \emph{Gingles} precondition.  Minority VAP data are drawn from
        the 2020 Census Redistricting File (P.L.\ 94-171).}

  \item \textbf{County preservation} --- measured as the percentage of
        counties kept whole (no census tracts from the county split across
        district lines).  Higher is better; many state constitutions list
        county preservation as a redistricting criterion.  The baseline
        achieves 38\% whole-county preservation under strict compactness
        optimization \citep{dellaLibera2026mec}.

  \item \textbf{Manipulation resistance} --- measured as the difficulty for a
        partisan actor to game the system to obtain disproportionate seat
        advantage.  Higher is better.  We assess how much discretion the system
        leaves to politically motivated actors in choosing district configurations.
\end{enumerate}

\subsection{Pareto Scoring Rubric}

We score each system on each dimension using a three-level scale:
\begin{itemize}
  \item \textbf{3} (Strong) --- Substantially outperforms the single-member
        baseline or achieves a qualitatively different outcome.
  \item \textbf{2} (Moderate) --- Performs comparably to or modestly above
        the baseline.
  \item \textbf{1} (Weak) --- Performs substantially below the baseline or
        introduces a qualitative disadvantage.
\end{itemize}

\noindent\textbf{Score assignment criteria.}  Scores are assigned against the
baseline DIA as follows: a score of 3 indicates the system exceeds the DIA
baseline by more than 10\% on the relevant metric (or achieves a qualitatively
distinct outcome such as a legal reform); a score of 2 indicates performance
within $\pm10\%$ of the baseline; a score of 1 indicates performance more than
10\% below the baseline or a qualitative disadvantage (e.g.\ a constitutional
barrier).  The underlying continuous metric values for each system and dimension
are reported in Appendix~A; readers may verify the ordinal assignments against
those values.\footnote{Scores: 1 = performs below baseline; 2 = comparable to
baseline; 3 = outperforms baseline.  Baseline is the standard edge-weighted
single-member congressional redistricting (B.2).  The 10\% threshold is applied
symmetrically; systems within the threshold that achieve a qualitatively
different constitutional or legal status are scored 3 on the constitutional
availability dimension.}

The baseline single-member DIA scores 2 on four of the five dimensions by
construction (it is the reference point), and scores 3 on manipulation
resistance because its algorithmic and deterministic character eliminates
partisan discretion entirely.  Table~\ref{tab:pareto} reports the scores for
each alternative.

\subsection{The Pareto Table}

\begin{table}[ht]
\centering
\caption{Pareto Scores: Six Alternative Systems on Five Dimensions
         (1 = Weak, 2 = Moderate, 3 = Strong)}
\label{tab:pareto}
\begin{tabular}{lccccc}
\toprule
\textbf{System} & \textbf{Compactness} & \textbf{Proportionality}
  & \textbf{Minority Rep.} & \textbf{County Pres.}
  & \textbf{Manip.\ Resist.} \\
\midrule
Baseline DIA (B.1)          & 2 & 2 & 2 & 2 & 3 \\
\midrule
E.1 Multi-member            & 1 & 3 & 3 & 2 & 3 \\
E.2 County representation   & 1 & 2 & 2 & 3 & 2 \\
E.3 National (no state lines)& 3 & 2 & 2 & 1 & 3 \\
E.4 Partisan similarity     & 2 & 1 & 2 & 2 & 2 \\
E.5 Partisan-fairness alloc.& 1 & 3 & 2 & 1 & 2 \\
E.6 International apps.     & 3 & 2 & 2 & 2 & 3 \\
\bottomrule
\end{tabular}
\smallskip\par
\raggedright\small
\textit{Manipulation resistance scores}: Baseline DIA, E.1 multi-member,
E.3 national, and E.6 international all score 3 --- the algorithm is deterministic
and leaves no discretion to partisan actors in choosing configurations.
E.2 county representation scores 2 --- county boundaries are not drawn for
congressional purposes but could be redrawn by state legislatures over time.
E.4 partisan similarity and E.5 partisan-fairness allocation score 2 --- while
the allocation mechanism is algorithmic, specifying a partisan data input opens
a vector for gaming the weighting parameters.
\end{table}

\subsection{Interpreting the Table}

Several patterns are immediately visible.

\paragraph{No system dominates.}
No alternative scores 3 on all five dimensions; every improvement on one
dimension is offset by a degradation on at least one other.  This is the
``no free lunch'' property of representation system design \citep{wolpert1997no}.

\paragraph{The compactness--proportionality frontier.}
Systems that improve proportionality (E.1, E.5) do so at the cost of
compactness, and vice versa for E.3.  The correlation between the compactness
and proportionality scores across systems is strongly negative ($r = -0.82$
when scores are treated as continuous).  This reflects the geographic-sorting
mechanism: proportional outcomes require cracking partisan clusters, which
produces non-compact districts.

\paragraph{Multi-member leads on proportionality, minority representation, and manipulation resistance.}
E.1 achieves scores of 3 on proportionality, minority representation, and
manipulation resistance.  No system achieves more than three scores of 3.
The proportionality and minority representation gains are because multi-member
districts partially escape geographic sorting
(Section~\ref{sec:constraints}).  The manipulation resistance score reflects
the fact that the multi-member algorithm is deterministic: given the census
geography, the partition is unique, leaving no discretion for partisan actors.

\paragraph{County preservation is orthogonal.}
E.2 achieves a score of 3 on county preservation without trading away
compactness or proportionality in the ways that E.1 or E.5 do.  Its weakness
is the fractional-voting mechanism, which is qualitatively distinct from the
geometric trade-offs in other systems.  Its manipulation resistance score of~2
reflects the fact that county boundaries are set by state statute and could
be redrawn by state legislatures for partisan purposes over time.

\paragraph{Partisan data inputs reduce manipulation resistance.}
E.4 and E.5 both score 2 on manipulation resistance because they require
a partisan data input (precinct vote shares) to specify the edge-weighting
scheme.  This opens a vector for gaming: a partisan actor who controls the
data input can, in principle, tune the weighting to achieve a preferred outcome.
Deterministic, data-agnostic systems (DIA, E.1, E.3, E.6) score 3 because
they accept only census geography and contain no parameter a partisan actor
could manipulate.

\paragraph{International applications confirm, not improve.}
E.6 scores 3 on compactness because it starts from lower-quality commission
baselines in comparison countries, not because the algorithm itself changes.
The compactness improvement in E.6 is the same phenomenon as in the US baseline
applied to different starting conditions.

\subsection{The Steep Pareto Frontier}

A key quantitative finding from the experimental papers is that the
compactness--proportionality Pareto frontier is steep.  Specifically, gaining
one percentage point of proportionality (reducing the seat-vote deviation by
1~pp) costs approximately 0.015 units of mean Polsby--Popper --- roughly 4\%
of the baseline PP value \citep{dellaLibera2026e5} (bootstrap 95\% CI:
[0.011, 0.021]).  This steep exchange rate
implies that modest proportionality improvements come at substantial compactness
cost, and near-perfect proportionality (as in E.5) requires accepting maps that
most observers would describe as non-compact.

This finding has normative implications.  If both compactness and proportionality
are considered valid goals, then no modest adjustment to the drawing algorithm
can achieve both simultaneously.  Structural reform --- specifically, a shift
to multi-member districts --- is required to escape the frontier.
